\documentclass[a4paper,12pt]{article}

\usepackage[T2A]{fontenc}
\usepackage[utf8]{inputenc}
\usepackage[russian]{babel}
\usepackage{amsmath,amssymb,mathtools}
\usepackage{helvet}
\usepackage{PTSans}
\renewcommand{\familydefault}{\sfdefault}
\usepackage{icomma}
\usepackage{geometry}
\geometry{left=20mm,right=20mm,top=20mm,bottom=22mm}
\usepackage[nopatch=footnote]{microtype}
\usepackage{tikz}
\usetikzlibrary{arrows.meta,positioning,shapes.geometric}
\usepackage{tabularx,booktabs}
\usepackage{enumitem}
\usepackage{hyperref}
\usepackage{parskip}
\usepackage{fancyhdr}
\usepackage{setspace}
\usepackage{xcolor}

\definecolor{accent}{HTML}{008080}
\definecolor{darkaccent}{HTML}{004D4D}
\definecolor{textgray}{HTML}{333333}
\definecolor{softgray}{HTML}{F1F4F4}

\hypersetup{colorlinks=true,linkcolor=darkaccent,urlcolor=darkaccent}
\onehalfspacing
\setlength{\parindent}{0pt}
\setlist[itemize]{leftmargin=1.6em,itemsep=0.25em,topsep=0.35em}
\setlist[enumerate]{leftmargin=1.8em,itemsep=0.35em,topsep=0.35em}
\renewcommand{\arraystretch}{1.35}

\setlength{\headheight}{25pt}
\pagestyle{fancy}
\fancyhf{}
\fancyhead[L]{\small\textcolor{textgray}{Кирилл Зиновьев --- 7 класс}}
\fancyhead[R]{\small\textcolor{textgray}{Математика}}
\fancyfoot[C]{\small\thepage}
\renewcommand{\headrulewidth}{0.4pt}
\renewcommand{\headrule}{\hbox to\headwidth{\color{accent}\leaders\hrule height \headrulewidth\hfill}}

\newcommand{\sectionline}{\vspace{-0.35em}\noindent\textcolor{accent}{\rule{\linewidth}{0.6pt}}\vspace{0.35em}}
\newcommand{\key}[1]{\textcolor{darkaccent}{\textbf{#1}}}

\tikzset{
  flowstart/.style={
    ellipse,
    draw=darkaccent,
    very thick,
    align=center,
    minimum width=34mm,
    minimum height=11mm,
    inner sep=4pt
  },
  flowproc/.style={
    rounded corners=3pt,
    rectangle,
    draw=accent,
    very thick,
    align=center,
    text width=86mm,
    minimum height=12mm,
    inner sep=5pt
  },
  flowprocsmall/.style={
    rounded corners=3pt,
    rectangle,
    draw=accent,
    very thick,
    align=center,
    text width=46mm,
    minimum height=12mm,
    inner sep=4pt
  },
  flowcond/.style={
    diamond,
    aspect=2.0,
    draw=accent,
    very thick,
    align=center,
    text width=42mm,
    inner sep=2pt
  },
  flowarrow/.style={
    -{Latex[length=3mm,width=2mm]},
    very thick,
    draw=darkaccent
  },
  flowlabel/.style={
    font=\small,
    fill=white,
    inner sep=1.5pt,
    text=darkaccent
  }
}

\begin{document}

\begin{titlepage}
\thispagestyle{empty}
\centering
\vspace*{18mm}

{\Large\textcolor{darkaccent}{\textbf{Математика, 7 класс}}}\par
\vspace{12mm}

{\Huge\textbf{Чек-лист}}\par
\vspace{5mm}

{\LARGE\textcolor{accent}{\textbf{Буквенные выражения, смысл формул и проценты}}}\par
\vspace{10mm}

{\large
Раскрытие скобок и приведение подобных слагаемых\par
Интерпретация буквенных выражений\par
Решение процентных задач пропорцией\par}

\vfill

{\large \today}\par
\vspace{10mm}

{\large\textbf{Лёвин Артём Александрович}}\par
{\large эксклюзивно для Кирилла Зиновьева}\par

\vspace{12mm}

\begin{tikzpicture}[x=1cm,y=1cm]
  \draw[accent,line width=1.2pt]
    (0,0)
    .. controls (3.0,0.8) and (7.5,-0.8) .. (11.0,0.1)
    .. controls (13.5,0.7) and (15.2,0.2) .. (16.5,-0.1);
\end{tikzpicture}

\end{titlepage}

\section*{1. Раскрытие скобок и подобные слагаемые}
\sectionline

\key{Подобные слагаемые} --- это слагаемые с одинаковой буквенной частью.
Складываются и вычитаются только их коэффициенты:
\[
6x-2x=4x,
\qquad
8y+3y+5y=16y.
\]

\key{Если перед скобками стоит «+»}, знаки внутри сохраняются:
\[
-3x+(7+4x+2)
=
-3x+7+4x+2
=
x+9.
\]

\key{Если перед скобками стоит «-»}, знак каждого слагаемого внутри скобок
меняется на противоположный:
\[
6x-(24+2x+10)
=
6x-24-2x-10
=
4x-34.
\]

Важно различать два действия:
\begin{itemize}
  \item сначала корректно раскрыть скобки;
  \item затем сгруппировать подобные слагаемые и выполнить вычисления
        с коэффициентами.
\end{itemize}

\subsection*{Быстрая проверка знаков}

Перед объединением подобных слагаемых задайте себе два вопроса:
\begin{enumerate}
  \item Был ли перед скобками внешний минус?
  \item Поменялся ли знак \emph{у каждого} слагаемого внутри этих скобок?
\end{enumerate}

\subsection*{Типичная ошибка}

Из записи
\[
-(2x+10)
\]
получается
\[
-2x-10,
\]
поскольку внешний минус действует на оба слагаемых.

\clearpage

\section*{Алгоритм: упростить выражение со скобками}

\begin{center}
\begin{tikzpicture}[node distance=8mm and 14mm]

  \node[flowstart] (start) {Начало};

  \node[flowproc,below=of start] (scan)
  {Найдите скобки и знак непосредственно перед ними};

  \node[flowcond,below=10mm of scan] (minus)
  {Перед скобками стоит «$-$»?};

  \node[
    flowprocsmall,
    below=18mm of minus,
    xshift=-50mm
  ] (change)
  {Поменяйте знак каждого слагаемого};

  \node[
    flowprocsmall,
    below=18mm of minus,
    xshift=50mm
  ] (keep)
  {Сохраните все знаки};

  \node[flowproc,below=46mm of minus] (group)
  {Соберите подобные слагаемые: отдельно буквенные, отдельно числовые};

  \node[flowproc,below=of group] (calc)
  {Выполните действия с коэффициентами и числами};

  \node[flowcond,below=10mm of calc] (check)
  {Все знаки и подобные слагаемые проверены?};

  \node[
    flowprocsmall,
    text width=39mm,
    left=8mm of check
  ] (fix)
  {Исправьте знаки и заново сгруппируйте};

  \node[flowstart,below=14mm of check] (finish)
  {Готово};

  \draw[flowarrow] (start) -- (scan);
  \draw[flowarrow] (scan) -- (minus);

  \draw[flowarrow]
    (minus)
    -|
    node[flowlabel,pos=0.25]{да}
    (change);

  \draw[flowarrow]
    (minus)
    -|
    node[flowlabel,pos=0.25]{нет}
    (keep);

  \draw[flowarrow] (change.south) |- (group.west);
  \draw[flowarrow] (keep.south) |- (group.east);

  \draw[flowarrow] (group) -- (calc);
  \draw[flowarrow] (calc) -- (check);

  \draw[flowarrow]
    (check)
    --
    node[flowlabel]{да}
    (finish);

  \draw[flowarrow]
    (check.west)
    --
    node[flowlabel,above]{нет}
    (fix.east);

  \draw[flowarrow] (fix.north) |- (calc.west);

\end{tikzpicture}
\end{center}

\clearpage

\section*{2. Видеть смысл за буквенным выражением}
\sectionline

Буквы обозначают величины. Поэтому полезно переводить запись
с математического языка на русский и обратно.

\subsection*{Прямоугольник со сторонами $a$ и $b$}

\begin{center}
\begin{tabularx}{\linewidth}{@{}>{\bfseries}l X@{}}
\toprule
Выражение & Смысл \\
\midrule
$a\cdot b$ & площадь прямоугольника \\
$2a+2b$ & периметр прямоугольника \\
$a+b$ & сумма длины и ширины \\
$2a$ & удвоенная длина \\
$a+2b$ & сумма длины и удвоенной ширины \\
$2a-b$ & разность удвоенной длины и ширины \\
\bottomrule
\end{tabularx}
\end{center}

\subsection*{Квадрат со стороной $x$}

\begin{center}
\begin{tabularx}{\linewidth}{@{}>{\bfseries}l X@{}}
\toprule
Выражение & Смысл \\
\midrule
$x^2$ & площадь квадрата \\
$4x$ & периметр квадрата \\
$2x$ & удвоенная сторона квадрата \\
$\dfrac{x}{2}$ & половина длины стороны квадрата \\
\bottomrule
\end{tabularx}
\end{center}

\subsection*{Треугольник со сторонами $x$, $y$, $z$}

\[
x+y+z=P,
\]
где $P$ --- периметр треугольника.

Запись
\[
2x+y+z
\]
уже не является периметром обычного треугольника со сторонами
$x$, $y$, $z$; это сумма удвоенной первой стороны и двух остальных сторон.

\subsection*{Контрольный вопрос}

Если выражение совпадает с известной формулой, назовите соответствующую
величину. Если готовой формулы нет, грамотно опишите выражение словами.

\clearpage

\section*{3. Проценты: три базовых типа задач}
\sectionline

\key{100\%} всегда соответствует целой величине, от которой считают проценты.

\subsection*{Тип 1. Найти процент от известного числа}

Пример: найти $40\%$ от $500$.

\[
\begin{array}{c|c}
100\% & 500\\
40\% & x
\end{array}
\qquad
\frac{x}{500}
=
\frac{40}{100}.
\]

Отсюда
\[
x=200.
\]

\subsection*{Тип 2. Найти число по известному проценту}

Пример: найти число, $6\%$ которого равны $54$.

\[
\begin{array}{c|c}
100\% & x\\
6\% & 54
\end{array}
\qquad
\frac{54}{x}
=
\frac{6}{100}.
\]

Следовательно,
\[
x
=
\frac{54\cdot100}{6}
=
900.
\]

\subsection*{Тип 3. Найти, сколько процентов одно число составляет от другого}

Пример: сколько процентов составляет $100$ от $500$?

\[
\begin{array}{c|c}
100\% & 500\\
x\% & 100
\end{array}
\qquad
\frac{x}{100}
=
\frac{100}{500}.
\]

Значит,
\[
x=20\%.
\]

\subsection*{Языковая подсказка}

Если в условии написано
«найдите $p\%$ \emph{от числа} $N$»,
то $N$ соответствует $100\%$.

Если написано
«найдите число, $p\%$ \emph{которого} равны $A$»,
искомое число соответствует $100\%$.

\clearpage

\section*{Алгоритм: решить задачу на проценты пропорцией}

\begin{center}
\begin{tikzpicture}[node distance=7mm and 14mm]

  \node[flowstart] (start)
  {Начало};

  \node[flowproc,below=of start] (read)
  {Определите целую величину, то есть то, чему соответствуют $100\%$};

  \node[flowcond,below=9mm of read] (knownwhole)
  {Целая величина известна?};

  \node[
    flowprocsmall,
    below=18mm of knownwhole,
    xshift=-50mm
  ] (setnum)
  {Поставьте известное число напротив $100\%$};

  \node[
    flowprocsmall,
    below=18mm of knownwhole,
    xshift=50mm
  ] (setx)
  {Поставьте $x$ напротив $100\%$};

  \node[flowproc,below=46mm of knownwhole] (second)
  {Заполните вторую пару «процент --- число»};

  \node[flowproc,below=of second] (prop)
  {Составьте пропорцию и перемножьте крест-накрест};

  \node[flowproc,below=of prop] (solve)
  {Найдите $x$; при возможности сократите вычисления};

  \node[flowcond,below=9mm of solve] (unit)
  {Ответ имеет требуемый смысл?};

  \node[
    flowprocsmall,
    text width=39mm,
    left=8mm of unit
  ] (fix)
  {Проверьте $100\%$ и пересоставьте пропорцию};

  \node[flowstart,below=14mm of unit] (finish)
  {Готово};

  \draw[flowarrow] (start) -- (read);
  \draw[flowarrow] (read) -- (knownwhole);

  \draw[flowarrow]
    (knownwhole)
    -|
    node[flowlabel,pos=0.25]{да}
    (setnum);

  \draw[flowarrow]
    (knownwhole)
    -|
    node[flowlabel,pos=0.25]{нет}
    (setx);

  \draw[flowarrow] (setnum.south) |- (second.west);
  \draw[flowarrow] (setx.south) |- (second.east);

  \draw[flowarrow] (second) -- (prop);
  \draw[flowarrow] (prop) -- (solve);
  \draw[flowarrow] (solve) -- (unit);

  \draw[flowarrow]
    (unit)
    --
    node[flowlabel]{да}
    (finish);

  \draw[flowarrow]
    (unit.west)
    --
    node[flowlabel,above]{нет}
    (fix.east);

  \draw[flowarrow] (fix.north) |- (prop.west);

\end{tikzpicture}
\end{center}

\clearpage

\section*{4. Итоговый чек-лист}
\sectionline

Перед тем как считать, проверьте:
\begin{itemize}
  \item внешний минус перед скобками меняет все внутренние знаки;
  \item внешний плюс перед скобками сохраняет знаки;
  \item подобные слагаемые имеют одинаковую буквенную часть;
  \item $ab$ для прямоугольника со сторонами $a$ и $b$ --- площадь;
  \item $2a+2b$ для прямоугольника --- периметр;
  \item $x^2$ для квадрата со стороной $x$ --- площадь;
  \item $4x$ для квадрата --- периметр;
  \item в процентной задаче сначала определяется величина,
        соответствующая $100\%$;
  \item ответ не подгоняется под готовый вариант:
        сначала выполняется собственное решение, затем проверка.
\end{itemize}

\subsection*{Самопроверка}

После решения полезно отдельно проверить:
\begin{enumerate}
  \item знаки после раскрытия скобок;
  \item арифметику при сложении положительных и отрицательных чисел;
  \item смысл полученного буквенного выражения;
  \item что в задаче на проценты правильно определено соответствие
        для $100\%$;
  \item единицу ответа:
        число, сантиметры, квадратные сантиметры или проценты.
\end{enumerate}

\clearpage

\section*{Тренировочный блок}
\sectionline

Упражнения расположены по нарастающей сложности.
Решения оформите письменно.

\begin{enumerate}

  \item Упростите выражение:
  \[
  7x-(18+3x+6).
  \]

  \item Упростите выражение:
  \[
  -4y+(9+7y-5).
  \]

  \item Приведите подобные слагаемые:
  \[
  12a-5b-7a+3b+4a.
  \]

  \item У прямоугольника длина равна $m$, ширина равна $n$.
  Объясните словами смысл каждого выражения:
  \[
  mn,
  \qquad
  2m+2n,
  \qquad
  m+2n.
  \]

  \item У квадрата сторона равна $x$.
  Какой геометрический смысл имеют выражения
  \[
  x^2,
  \qquad
  4x,
  \qquad
  \frac{x}{2}?
  \]

  \item Найдите $35\%$ от числа $480$.

  \item Найдите число, $12\%$ которого равны $72$.

  \item Сколько процентов составляет число $45$ от числа $180$?

  \item Упростите выражение и затем при $x=4$ найдите его значение:
  \[
  3(2x-5)-(x+7)+2(x-1).
  \]

  \item У прямоугольника длина $a$ см, ширина $b$ см.
  Известно, что
  \[
  a=12,
  \qquad
  b=7.
  \]
  Найдите значение выражения $2a+2b$ и объясните его геометрический смысл.
  Затем определите, сколько процентов ширина составляет от длины.

\end{enumerate}

\clearpage

\section*{Ответы к тренировочному блоку}
\sectionline

\begin{center}
\begin{tabularx}{\linewidth}{@{}c X@{}}
\toprule
\textbf{№} & \textbf{Ответ} \\
\midrule

1 &
$4x-24$
\\

2 &
$3y+4$
\\

3 &
$9a-2b$
\\

4 &
$mn$ --- площадь прямоугольника;
$2m+2n$ --- периметр;
$m+2n$ --- сумма длины и удвоенной ширины.
\\

5 &
$x^2$ --- площадь квадрата;
$4x$ --- периметр;
$x/2$ --- половина длины стороны.
\\

6 &
$168$
\\

7 &
$600$
\\

8 &
$25\%$
\\

9 &
После упрощения:
$7x-24$;
при $x=4$:
$4$.
\\

10 &
$2a+2b=38$ см --- периметр прямоугольника;
ширина составляет
\[
\frac{7}{12}\cdot100\%
=
58\frac{1}{3}\%.
\]
\\

\bottomrule
\end{tabularx}
\end{center}

\end{document}